% ══════════════════════════════════════════════════════════════════════════════
% CHAPTER 2 — LITERATURE REVIEW (FINAL CLEAN VERSION)
% ══════════════════════════════════════════════════════════════════════════════
\chapter{Literature Review}

\section{MSME Landscape and Statistical Analysis}

The MSME sector is a very key component of India's economy. It puts out around 30.1\% of the national GDP which is a very large number. Also it accounts for almost 35.4\% of what we see in manufacturing output and also about 45.73\% of total exports. We have in India around 6.5 crore registered MSME units and these businesses together put to work almost 28 crore people as of \cite{pib_msme2025}. In terms of scale and impact it is safe to say that this sector is very at the core of the Indian economy.

One to note is the rate at which formalization is taking off. By April 2026 we saw registration of about 7.94 crore enterprises via the Udyam Registration and the Udyam Support Platform (UAP). That is a large scale increase and it is a sign that we are seeing greater recognition and integration of small scale businesses into the main economic structure \cite{msme_annual_report}. \begin{table}[htbp].

\centering

\small

\begin{tabular}{|l|r|}

\hline

\textbf{Category} & \textbf{Number of Enterprises} \\

\hline

Total MSMEs & 7,94,25,711 \\

Udyam Registered & 4,72,77,069 \\

UAP Registered & 3,21,48,642 \\

\hline

\end{tabular}

\caption{MSME Registration Statistics \cite{msme_annual_report}}

\end{table}

If you take a closer look at the distribution you'll see what I mean in that field we have almost only micro enterprises.


\begin{table}[htbp]
\centering
\small
\begin{tabular}{|l|r|}
\hline
\textbf{Enterprise Type} & \textbf{Count} \\
\hline
Micro Enterprises & 7,88,97,394 \\
Small Enterprises & 4,91,260 \\
Medium Enterprises & 37,057 \\
\hline
\end{tabular}
\caption{MSME Distribution by Enterprise Category \cite{msme_annual_report}}
\end{table}

What we see is that by and large the preponderance of businesses in India are very small. They work with limited resources and also for the most part do not have much in terms of technology infrastructure. Although we see high adoption of smart phones by business owners today what they do with technology is still very basic. It is mostly for communication and may also include digital payments not for the running of the business like order tracking, customer management, or inventory. In large part those functions are still performed manually.

In the MSME sector we see that it is micro and small businesses which use what they are familiar with, what is everyday to them. There is a true need for basic digital solutions which may help them to work better without at the same time requiring them to become tech savvy over night.

% ─────────────────────────────────────────────────────────────────────────────

\section{Informal Business Characteristics}

In India a large number of MSMEs operate in the informal sector which include small scale and home based businesses. These businesses are mostly run by a single individual or a family, have low investment, a small work force and do not have formal digital systems. Because of their informal structure these companies use mainly manual processes which include hand written notes and direct messages for order and booking management. In this segment technology adoption is very much dependent on how simple and familiar the system is. Companies do not adopt complex set up or technical tools which require great IT skill, instead they prefer tools that fit into their present work flow.

% ─────────────────────────────────────────────────────────────────────────────

\section{Adoption of WhatsApp in Small Businesses}

WhatsApp has become a global leader in terms of which communication platform is used, with over 2 billion monthly active users. In India it is the largest digital communication player, with a reported user base of over 530 million which we may assume to be nearly all smartphone users there \cite{hyperleap_india2026}. We see that daily on the platform there are over 100 billion messages sent out which speaks to the scale and reliability of the platform \cite{infobip_stats2026}. Also in the area of business communication we see very high adoption. Over 200 million businesses world wide use WhatsApp Business accounts and also a large number of small and medium sized businesses use it for every day customer communication \cite{infobip_stats2026}. Also we see that over 175 million users are interacting with businesses daily on WhatsApp which plays a large role in customer engagement \cite{d7_stats2026}.

WhatsApp reports to have very high engagement numbers. We see open rates of up to 98\% for business messages and also that they get much better response rates which is a great improvement over what we see in traditional communication methods like email and SMS.

\begin{table}[htbp]
    \centering
    \begin{tabular}{lccc}
        \hline
        \textbf{Metric} & \textbf{WhatsApp} & \textbf{Email} & \textbf{SMS} \\
        \hline
        Open Rate & 98\% & 22\% & 95\% \\
        Response Rate & 45\% & 6\% & 12\% \\
        Click-through Rate & 35\% & 3\% & 8\% \\
        Conversion Rate & 15\% & 2\% & 4\% \\
        \hline
    \end{tabular}
    \vspace{0.2cm}
    \caption{Engagement Comparison: WhatsApp vs. Traditional Channels \cite{hyperleap_india2026}}
    \label{tab:engagement_stats}
\end{table}

These engagement numbers prove that which WhatsApp has become the go to platform for business communication. We see that our users are on the app daily which in turn makes it very effective for real time communication and order management.

% ─────────────────────────────────────────────────────────────────────────────

\section{Challenges Identified in Existing Practices}

Despite wide scale adoption most MSMEs are still using WhatsApp in a very unstructured and manual way. Business owners use the single chat interface for customer service, order info, and payments which in turn causes issues and increases their work load. Also we see that customer expectations have grown. Immediacy is what makes the platform successful we see that the large majority of messages are read very quickly. We also note that consumers expect prompt response; many expect reply within an hour and they will take their business elsewhere if they don't get it which is reported in \cite{hyperleap_india2026}. This puts great pressure on small business owners who are trying to manage many conversations at once without support of automation.

While in the case of WhatsApp Business API which provides chatbots, multi agent support, and CRM integration they do so at a high cost and also from a technical point of view it is a complex solution. Also these platforms require business to go through a verification process, API integration, and system configuration which in turn makes them a no go for what we term as informal and small scale businesses. Thus we see that adoption of what may be classified as advanced automation features is very low in micro businesses.

This creates a clear gap use of WhatsApp and lack of accessible automation tools for MSMEs. Addressing this gap forms the foundation for the proposed ChatServe system.